\documentclass{article}
\usepackage[UTF8]{ctex}
\usepackage{amsmath,amssymb,amsfonts}
\usepackage{enumitem}
\usepackage{geometry}
\geometry{a4paper, margin=1in}
% === 错题本排版优化（不参与题目标识） ===
\setlength{\parindent}{0pt}
\setlength{\parskip}{0.35em plus 0.10em minus 0.08em}
\setlength{\emergencystretch}{3em}
\setlist{itemsep=0.25em, topsep=0.40em, parsep=0pt}
\allowdisplaybreaks[2]
% === 错题本排版优化结束 ===
\begin{document}
\section*{天津大学2024年数学分析真题}
\section*{58．天津大学2024年数学分析试卷}

\subsection{题目1}
1．构造一个二元函数满足以下条件并证明．

\begin{itemize}
  \item[（1）] 在原点处连续且偏导数存在；
  \item[（2）] 在原点处不可微．
\end{itemize}

\subsection{题目2}
2．设 \(f(x)\) 在 \([a, b]\) 上二阶可导，\(f(a)=f(b)\) ，且 \(\left|f^{\prime \prime}(x)\right| \leq c\) ，证明：\(\left|f^{\prime}(x)\right| \leq \frac{1}{2} c(b-a)\) ．

\subsection{题目3}
3．设 \(f(x)\) 在 \([a,+\infty)\) 上可导，\(f(a)>a\) ，且 \(f^{\prime}(x)+f(x)>1+x\) ，证明：\(f(x)>x\) ．

\subsection{题目4}
4．设 \(f(x)\) 在 \([1,+\infty)\) 上一致连续，证明：存在 \(M>0\) ，满足 \(\frac{|f(x)|}{x} \leq M, x \in[1,+\infty)\) ．

\subsection{题目5}
5．设数列 \(\left\{a_{n}\right\}\) 满足 \(\lim _{n \rightarrow \infty} \frac{a_{n}}{n^{\alpha}}=1(\alpha>0)\) ，求极限 \(\lim _{n \rightarrow \infty} \frac{a_{1}+a_{2}+\cdots+a_{n}}{n^{1+\alpha}}\) ．

\subsection{题目6}
6．设 \(f(x)\) 是以 \(2 \pi\) 为周期且在 \([-\pi, \pi]\) 上可积的函数，\(a_{n}, b_{n}\) 为 \(f(x)\) 的傅里叶级数．

\begin{itemize}
  \item[（1）] 求 \(f(x+1)\) 的傅里叶系数．
  \item[（2）] 求 \(f(2 x)\) 的傅里叶系数．
\end{itemize}

\begin{itemize}

\subsection{题目7}
\end{itemize}

\subsection{题目7}
原题边界“7．”：设 \(f(x)\) 为 \([a, b]\) 上的单调函数，\(g(x)\) 是周期为 \(T\) 的连续函数，且 \(\int_{0}^{T} g(x) \mathrm{d} x=0\) ，证明：
\[
\lim _{\lambda \rightarrow+\infty} \int_{a}^{b} f(x) g(\lambda x) \mathrm{d} x=0
\]

\subsection{题目8}

\subsection{题目8}
原题边界“8．”：已知序列 \(\left\{x_{n}\right\},\left\{y_{n}\right\}\) 满足 \(x_{n+1}=y_{n}+\frac{1}{2} x_{n}\) ，证明：\(\left\{y_{n}\right\}\) 收敛的充要条件是 \(\left\{x_{n}\right\}\) 收敛．




\subsection{题目9}

\subsection{题目9}
原题边界“9．”：记 \(D_{R}=\{|x| \leq R,|y| \leq R,|z| \leq R \mid R>0\}\) ．求极限
\[
\lim _{R \rightarrow+\infty} \iiint_{D_{R}}(x+y+z)^{2} e^{-\left(x^{2}+y^{2}+z^{2}\right)} \mathrm{d} x \mathrm{~d} y \mathrm{~d} z
\]

\subsection{题目10}

\subsection{题目10}
原题边界“10．”：求积分 \(\int_{0}^{+\infty} \frac{(\sin x)^{2}}{x^{2}} \mathrm{~d} x\) ．

\subsection{题目11}
11．设 \(L\) 为 \(x^{2}+y^{2}+z^{2}=2 R x\) 与 \(x^{2}+y^{2}=2 r x\) 的交线 \((R>r>0)\) 的上半部分 \((z \geq 0)\) ，求曲线积分

\[
\int_{L}\left(y^{2}+z^{2}\right) \mathrm{d} x+\left(z^{2}+x^{2}\right) \mathrm{d} y+\left(x^{2}+y^{2}\right) \mathrm{d} z .
\]

从 \(z\) 轴正向看，\(L\) 取逆时针方向。

\subsection{题目12}
12．设 \(\left\{f_{n}(x)\right\}\) 为定义在 \([a, b] \backslash\left\{x_{0}\right\}\) 上的函数列，且在 \([a, b] \backslash\left\{x_{0}\right\}\) 上一致收敛于 \(f(x)\) ，同时

\[
\lim _{x \rightarrow x_{0}} f_{n}(x)=c_{n}(n=1,2, \cdots) .
\]

（1）证明： \(\lim _{n \rightarrow \infty} c_{n}\) 存在．\\
（2）证明： \(\lim _{x \rightarrow x_{0}} f(x)\) 存在，且 \(\lim _{x \rightarrow x_{0}} f(x)=\lim _{n \rightarrow \infty} c_{n}\) ．
\end{document}
